\scanpage{37}
\mspage{29}
% L1–2
via le diagramme $\Delta$ et l'ens $V^*$ de cardinal 3.

% L3–7
\uline{Lemme 2} \quad a) $\forall x\in C-\Delta$, soit $\Delta(x)=\Delta\cap E(x)_0$. Alors $[\operatorname{card}\Delta(x)=3,\ \text{et}]$ les deux applications $\prod_{j\in I}T_j\simeq\Delta\hookrightarrow\!\longrightarrow T_i$ $(i\in I)$ définissent des bijections $\Delta(x)\simeq T_i$ i.e. $\Delta_x$ est le graphe d'une bijection $T_i\simeq T_{i'}$ (si \deleted{\unsure{$i,i'$|$i\ne i'$}}, $I=\{i,i'\}$).
% L8–14
b) Soit \deleted{\unsure{l'ens}} $\deleted{\unsure{E}}\unsure{F'|T'|\mathbb F'}$ l'ens des parties \inserted{$\Delta'$} de $\Delta$ \deleted{\unsure{$(x)$}} ayant les propriétés précédentes [i.e. de cardinal 3, et formé de p$^{\rm ts}$ 2 à deux non liés \inserted{ds $\Delta$}], et \unsure{considérer|considérons} les applications canoniques $(C-\Delta)\longrightarrow F$ et $(C-\Delta)\longrightarrow V^*$, d'où $(C-\Delta)\longrightarrow F'\times V^*$. Cette application est bijective.

% L15–22
Il faut expliciter la relation de \deleted{\unsure{incidence}} \inserted{liaison} dans $C\simeq\Delta\amalg(F'\times V^*)$. On a déjà la relation de liaison pour deux él \deleted{\unsure{qui}} \inserted{$x,y$ si $x,y$} sont dans $\Delta$, ou si $x\in\Delta$, $y\unsure{\in|=}(\Delta(y),v(y))\in F'\times V^*$ [c'est la relation $x\in\Delta(y)$]. Il reste à l'expliciter si $x,y\in F'\times V^{\unsure{\wedge|*}}$, $x=(\Delta(\deleted{\unsure{y}}x),v(x))$, $y=(\Delta(y),v(y))$. On trouve que $x,y$ \inserted{distincts} sont liés ssi
\[
\left\{\begin{aligned}
&\text{ou bien }\Delta(x)\cap\Delta(y)\ne\varnothing,\quad v(x)=v(y)
&&(\text{donc }\operatorname{card}\Delta(x)\cap\Delta(y)=1)\\
&\text{ou bien }\Delta(x)\cap\Delta(y)=\varnothing,\quad v(x)\ne v(y)
\end{aligned}\right.
\]
% archive sheet number: 36
